\manualchapter{Noise and clock reference}{sec:registers}
\begin{table}[ht]\centering
\begin{tabularx}{\textwidth}{@{}lX@{}}
\toprule Noise setting & Pattern selection \\\midrule
0 & 5-bit then long polynomial \\
1, 3 & 5-bit polynomial \\
2 & 5-bit then 4-bit polynomial \\
4 & Long polynomial \\
5, 7 & Pure tone \\
6 & 4-bit polynomial \\\bottomrule
\end{tabularx}
\caption{Noise-register combinations. The long polynomial is 17-bit by
default or 9-bit with the global LFSR switch.}
\label{tab:shift-register}
\end{table}

\begin{table}[ht]\centering
\begin{tabularx}{\textwidth}{@{}lX@{}}
\toprule Global switch & Effect when enabled \\\midrule
Low Frequency & Selects the slower shared divider. \\
High-Pass 2 from 4 & Uses Tone 4 to filter Tone 2. \\
High-Pass 1 from 3 & Uses Tone 3 to filter Tone 1. \\
Tone 3 High Frequency & Runs Tone 3 from the faster clock. \\
Tone 1 High Frequency & Runs Tone 1 from the faster clock. \\
LFSR & Selects the shorter 9-bit long polynomial. \\\bottomrule
\end{tabularx}
\caption{Exposed audio-control register switches.}
\label{tab:audio-control}
\end{table}

Choose a noise setting with the global switches unchanged, then vary one
switch at a time. The switches affect the synthesis clock and coupling;
they are not independent post-processing effects. Settings 5 and 7 are
useful references when comparing a noisy patch with a simple tone.
